% Generated by scripts/plot.py -- do not edit, run `make figures`.
% Machines found in results/: LAP1184, apple-m5
\expandafter\gdef\csname PC@cross@tabcap@machines\endcsname{%
Cấu hình các máy dùng để đo. Thông tin do \texttt{scripts/machine\_info.sh} ghi lại lúc chạy đợt đo; kích thước cache line do hệ điều hành báo.%
}
\expandafter\gdef\csname PC@cross@tab@machines\endcsname{%
\begingroup\footnotesize
\begin{tabular}{@{}>{\raggedright\arraybackslash}p{4.4cm}>{\raggedright\arraybackslash}p{4.95cm}>{\raggedright\arraybackslash}p{4.95cm}@{}}
\toprule
 & \textbf{\texttt{LAP1184}} & \textbf{\texttt{apple-m5}} \\
\midrule
CPU & 13th Gen Intel(R) Core(TM) i5-1335U & Apple M5 \\
Kiến trúc & x86\_64 & arm64 \\
Số core & 12 luồng phần cứng trên 6 core vật lý & 10 (4 P-core + 6 E-core) \\
Cache line (hệ thống báo) & 64 B & 128 B \\
Cache line dùng để đệm trong mã & 64 B & 128 B \\
Cache L1d & 48K & 128 KiB (P) / 64 KiB (E) \\
Cache L2 & 1280K & 16 MiB per 4 P cores / 6 MiB per 6 E cores \\
Cache L3 & 12288K & – \\
RAM & 7.607 GiB & 16 GiB \\
Hệ điều hành & Ubuntu 20.04 LTS & macOS 27.0.1 (26A434) \\
Nhân & Linux 5.15.153.1-microsoft-standard-WSL2 & Darwin 27.0.0 \\
Compiler & c++ (Ubuntu 11.5.0-10ubuntu1\textasciitilde{}20\textasciitilde{}ppa3) 11.5.0 & Apple clang version 21.0.0 (clang-2100.1.1.101) \\
Cờ biên dịch & \texttt{-std=c++20 -O2 -Wall -Wextra -pthread -fopenmp} & \texttt{-std=c++20 -O2 -Wall -Wextra -pthread} \\
OpenMP & có & không (không có baseline OpenMP) \\
Lệnh chờ trong vòng spin & \texttt{pause} & \texttt{isb} \\
Ghim thread vào core & có (\texttt{sched\_setaffinity}) nhưng các benchmark không dùng & không có (macOS trên Apple Silicon không có API ghim thread) \\
Nguồn điện lúc đo & cắm sạc & cắm sạc \\
Thành viên chạy đo & \PCowner{LAP1184} & \PCowner{apple-m5} \\
\bottomrule
\end{tabular}
\endgroup%
}
\expandafter\gdef\csname PC@cross@figcap@lock\endcsname{%
So sánh giữa các máy, phần lock: thời gian đo chia cho cận dưới lý tưởng của chính máy đó, \texttt{OUTSIDE\_WORK} = 0, mỗi ô một máy, cùng thang đo. Chia cho cận dưới khử khác biệt về tốc độ core, nên các ô so được với nhau dù số vòng lặp được hiệu chỉnh riêng cho từng máy. Vạch dọc liền: số core của từng máy; vạch chấm: số P-core hoặc số core vật lý. Điểm có vòng tròn bao quanh: cấu hình sụp đổ (hoặc đã vượt thời hạn ở số vòng lặp đầy đủ) nên được đo với số vòng lặp nhỏ hơn 100 lần; giá trị vẽ đã quy đổi tuyến tính về cùng số vòng lặp với các điểm khác.%
}
\expandafter\gdef\csname PC@cross@figcap@barrier\endcsname{%
So sánh giữa các máy, phần barrier: số episode mỗi giây theo số thread, mỗi ô một máy, cùng thang đo. Vạch dọc liền: số core của từng máy; vạch chấm: số P-core hoặc số core vật lý.%
}
\expandafter\gdef\csname PC@cross@figcap@counter\endcsname{%
So sánh giữa các máy, phần counter ở mức tranh chấp cao nhất (\texttt{OUTSIDE\_WORK} nhỏ nhất của từng máy): thời gian chạy quy về $10^6$ lần tăng mỗi thread, mỗi ô một máy, cùng thang đo. Vạch dọc liền: số core của từng máy; vạch chấm: số P-core hoặc số core vật lý.%
}
\expandafter\gdef\csname PC@cross@val@hosts\endcsname{%
LAP1184,apple-m5%
}
\expandafter\gdef\csname PC@cross@val@nhosts\endcsname{%
2%
}
